\section{Methods}

\begin{sidewaysfigure}\centering
  \includegraphics[width=0.98\textheight]{conceptual_model.pdf}
  \caption{Conceptual model of the coupled SWAT\textsuperscript{+}/MODFLOW~6 PFAS
  fate-and-transport framework. The surface-water leg partitions PFAS in the soil
  profile across the aqueous, solid, and air--water-interface phases and routes it
  to streams with runoff, lateral flow, and sediment. New in this work
  (highlighted): a process-based PFAS engine ported to SWAT\textsuperscript{+};
  PFAS groundwater transport in MODFLOW~6 (replacing RT3D); and the first two-way
  SWAT\textsuperscript{+}--MODFLOW~6 coupling --- the daily recharge/baseflow water
  exchange, the retardation-lagged vadose pathway, the legacy groundwater plume,
  and the groundwater PFAS discharge to streams.}
  \label{fig:conceptual}
\end{sidewaysfigure}

\subsection{The SWAT\textsuperscript{+} PFAS surface-water engine}
\label{sec:engine}

\subsubsection{Design principle}
The engine is implemented in modern, free-form Fortran SWAT\textsuperscript{+} as a constituent that rides the existing spatial-object and constituent-hydrograph infrastructure. Its design separates the two halves of the problem along the line of where PFAS behavior is genuinely distinctive. The soil column---where PFAS partition simultaneously among water, solids, and the air--water interface---implements the process-based three-phase model of \citet{rafiei2023pfas}, ported from legacy SWAT (where it was first developed within a SWAT--MODFLOW--RT3D framework) into the modern SWAT\textsuperscript{+} engine. The channel and reservoir network---where PFAS behave, to first order, as a strongly sorbing conservative tracer with settling and benthic exchange---reuses the mathematics of the established SWAT\textsuperscript{+} pesticide subsystem, restricted to its linear-partition subset. This separation makes the engine both economical (most of the routing machinery already exists and is tested) and auditable (the ported code is confined to a small, independently verifiable set of routines).

\subsubsection{Soil three-phase equilibrium}
For each hydrologic response unit (HRU) and soil layer, the total PFAS mass present at the start of a day is partitioned among three coexisting phases at local equilibrium: a dissolved aqueous phase, a solid-sorbed phase described by a Freundlich isotherm, and an interfacially adsorbed phase at the air--water interface described by a Langmuir isotherm, the air--water interfacial term being the mechanism that can dominate PFAS retention in unsaturated media \citep{brusseau2018,guo2020,rafiei2023pfas}. Writing $C_w$ for the aqueous concentration, the layer mass balance is
\begin{equation}
M \;=\; \theta\,V\,C_w \;+\; m_s\,k_f\,C_w^{\,n} \;+\; \Gamma_{\max}\,\frac{K_L\,C_w}{1+K_L\,C_w}\,A_{aw},
\label{eq:threephase}
\end{equation}
where $M$ is the layer PFAS mass, $\theta$ the volumetric water content, $V$ the layer volume, $m_s$ the solid mass, $(k_f,n)$ the Freundlich coefficient and exponent, $(\Gamma_{\max},K_L)$ the Langmuir maximum surface excess and affinity, and $A_{aw}$ the air--water interfacial area, estimated from soil water content and median grain diameter. Because the right-hand side of Eq.~\eqref{eq:threephase} is strictly monotonic increasing in $C_w>0$, the equation has a unique positive root, which the engine isolates by a bracketed bisection and accelerates with an Improved-Halley iteration, retaining the bisection bracket as a guaranteed fallback. The original SWAT-MODFLOW-RT3D implementation solved this balance in quadruple precision \citep{rafiei2023pfas}; we re-implemented it in double precision and verified that the choice is safe: against a 40-digit decimal reference, the double-precision solver holds a worst-case relative error of $1.7\times10^{-7}$ across $29{,}160$ cases spanning the stiff corners of the parameter space (Freundlich $n$ down to $0.31$, layer mass down to $10^{-6}$~kg~ha\textsuperscript{-1}, and dry-to-saturated water contents), with no non-convergence. Once $C_w$ is known, the dissolved, solid, and interfacial masses follow by back-substitution and set the concentrations available to each loss pathway.

\subsubsection{Land-phase mobilization}
Following the equilibrium solve, dissolved PFAS is removed from the surface layer in proportion to surface runoff and percolation through a runoff/leaching partition coefficient, from deeper layers in proportion to lateral subsurface flow, and the solid-sorbed fraction is removed in proportion to sediment yield by an enrichment-ratio formulation, exactly paralleling the SWAT\textsuperscript{+} pesticide land phase so that the soil pool decrements consistently across all four pathways. The four daily losses---surface runoff, lateral flow, leaching, and sediment-bound---are accumulated per HRU and per PFAS compound and become the constituent loads delivered to the channel network. A whole-run land-phase mass balance confirms closure: for every HRU the change in soil pool equals the sum of the four pathway losses to a relative residual of order $10^{-5}$.

\subsubsection{In-stream transport and cross-validation}
In the channel, PFAS partitions between a dissolved and a sediment-sorbed fraction by a linear organic-carbon coefficient, $f_{\text{sol}}=1/(1+k_d\,c_s)$, where $k_d$ is the partition coefficient scaled by reach bed-carbon fraction and $c_s$ the suspended-sediment concentration, following the organic-carbon normalization of sediment--water partitioning for hydrophobic pollutants \citep{karickhoff1979,chapra1997}. The sorbed fraction settles to the bed; bed PFAS resuspends, diffuses across the sediment--water interface, and is permanently lost by burial; the dissolved and sorbed fractions leave the reach in the outflow. These are precisely the non-reactive terms of the SWAT\textsuperscript{+} pesticide channel solver, with the pesticide-specific aqueous decay, volatilization, and metabolite production removed because PFOS is non-volatile and recalcitrant on the relevant timescale. We exploit this correspondence as a correctness check. Configuring a pesticide with unit decay factors, zero volatilization, and no metabolites reduces the pesticide solver exactly to the PFAS linear-partition subset; driving both routines with identical inflow and parameters over a 60-day synthetic reach sweep, including a dry-reach day, the new PFAS routine and the established pesticide routine agree on all nine shared mass terms to a maximum absolute difference of $2.4\times10^{-7}$~kg and a maximum relative difference of $2.6\times10^{-7}$---single-precision floating-point tolerance. The in-stream mathematics is therefore not novel and not a source of error; the novelty is the soil column that supplies its loads.

\subsubsection{Reservoirs and point sources}
Reservoirs are treated as mass-conserving completely mixed cells: the PFAS entering a reservoir on a given day is added to the stored mass and the outflow carries the fraction leaving with the day's discharge, with the remainder retained, so that no mass is created or destroyed at impoundments (a benthic burial sink can be enabled but is omitted here, PFOS being mobile). Point sources---contaminated-site leachate, wastewater effluent, or firefighting-foam discharge---are injected as a per-reach daily load read from an optional input file and added to the reach inflow before routing; the injected mass is tracked in the in-stream mass balance. Because routing PFAS through a real network requires the constituent-hydrograph infrastructure that SWAT\textsuperscript{+} allocates only when a constituent is declared, the engine activates that infrastructure itself when PFAS is present, so a PFAS simulation requires no dummy constituent. A whole-network in-stream balance closes: cumulative reach inflow minus outflow equals the sum of permanent burial and the change in channel storage to floating-point tolerance, and the point-source input line reconciles exactly with the prescribed load. Daily per-reach output reports the PFAS mass balance and the outflow concentration in ng~L\textsuperscript{-1}, the quantity compared against ambient observations. This process-based surface-water engine supplies the land-derived in-stream PFAS load; the groundwater pathway developed below adds the legacy-plume and baseflow-discharge contribution to the same reaches.
